\chapter{Methodology}

This project studies arrest outcomes following NYPD vehicle stops in New York City. We study \textbf{whether the observed stop resulted in an arrest and which characteristics help explain or predict that outcome}. The statistical target is therefore $P(\texttt{arrested}=1 \mid X)$, where $X$ contains the predictor's information.

\paragraph{Data Origin:} The data comes from the \textit{NYPD Vehicle Stop Reports}, published by the City of New York. The raw file used for the project contains 2,488,293 stop records and 19 attributes covering time, location, officer command, vehicle type, motorist demographics and enforcement outcomes.

\paragraph{Data Context and Problems:} The data covers only about three years, it provides no complete description of the reason for the stop, and it contains partially missing demographic information. Coordinates are also incomplete or rounded in a non negligible fraction of observations. 

\paragraph{Data Preparation:} The main objective was to retain information that could help us generalize arrest-related variables, while removing identifiers, overfitting predictors and features that could increase model complexity without adding interpretable signal. Here, we sum up what we did:

\begin{center}
\begin{tabular}{lll}
\hline
\textbf{Block} & \textbf{Main operation} & \textbf{Final representation} \\
\hline
Target  & retain arrest only & $\texttt{arrested}\in\{0,1\}$ \\
Time & parse and decompose & year, month, weekday, hour \\
Demographics & clean and group sparse values & age, age bracket, sex, ethnicity \\
Geography & aggregate high cardinality command & area, boro, coordinates for EDA \\
Vehicle & remove commercial and TLC groups & 2 wheels, 4 wheels \\
\hline
\end{tabular}
\end{center}

In particular, we must mention that missing demographic values were preserved through explicit Unknown levels (when appropriate), because we considered that the complete deletion would alter the empirical population instead of removing noise. Furthermore, the bracketing of commands was done based on a report published by the NYPD, and the categories of vehicles that we retain are either any vehicles with 2 wheels or private vehicles with 4 wheels. These two macro categories originated 2 separate datasets.
After cleaning, we understood that the data available was very difficult to deal with. Most useful predictors are categorical and the main numerical variables, age and hour, are not naturally represented by a simple linear effect. At that point, we chose to go for parsimonious and interpretable models.

\paragraph{Splitting the Analysis:} The cleaned datasets (from the data preparation) are \textbf{analysed as two separate populations with the same response and the same research focus}. This split is mainly justified by the strong differences in size, arrest association, age composition and temporal/geographical distributions.

\paragraph{Exploratory Data Analysis and Goal:} The EDA is used to identify the structure that should be carried into the models. Its main purpose is to determine which predictors deserve attention, where the two datasets agree or disagree, and which relationships do not show linearity.

\paragraph{Univariate Analysis:} First of all we have huge difference in scale and class prevalence. The two wheeled dataset contains 189,610 stops, with 10,560 arrests, corresponding to approximately $5.6\%$. The four wheeled dataset contains 2,157,445 stops, with 70,600 arrests, approximately $3.3\%$. \textbf{Arrest is, thus, a rare outcome in both populations} and especially rare for four wheeled vehicles. Secondly, the distribution of the predictors also highlight how the two dataset differ. \textbf{Two wheeled stops are concentrated among younger motorists} and exhibit pronounced seasonal variation, with much higher volumes during warmer months and strong growth from 2023 to 2025. \textbf{Four wheeled stops have a broader age distribution and a much more stable annual and seasonal profile}. Geographic distribution of stops is also more even for four wheeled vehicles.
These differences justify the separate analyses but, at the same time, both datasets share the same difficulty: a rare binary target must be explained using mostly categorical or discrete predictors.

\paragraph{Bivariate Analysis:} We first use Cramer's $V$ to summarize pairwise association among categorical predictors. For two wheeled vehicles, the clearest pairwise associations with arrest are geographical location ($V\approx0.13$ for area and boro), hour or time bracket ($V\approx0.11$) and ethnicity ($V\approx0.10$). Sex ($V\approx0.02$) and age or age bracket ($V\approx0.04$) are weak when considered alone. For four wheeled vehicles, ethnicity remains around $V\approx0.10$, hour around $0.09$, age bracket around $0.07$, sex around $0.05$ and geography around $0.05$ to $0.06$. The matrices also show that \textbf{predictor to predictor associations are generally modest}. The largest recurring structure is between ethnicity and geography, approximately $0.17$ to $0.19$. Pearson correlations among the raw numerical variables are similarly weak. This is useful for diagnostics purposes but it also suggests that simple linear target-predictor relationships do not suffice to explain the target. The bivariate results therefore focus the remainder of the analysis on five axes: $\texttt{boro}$, $\texttt{hour}$, $\texttt{sex}$, $\texttt{ethnicity}$ and $\texttt{age}$ or $\texttt{age\_bracket}$.

\begin{figure}[H]
\centering
\begin{minipage}{0.49\textwidth}
\centering
\includegraphics[width=\textwidth]{img/categ_corr_2w.png}
\small (a) Two wheeled vehicles
\end{minipage}
\hfill
\begin{minipage}{0.49\textwidth}
\centering
\includegraphics[width=\textwidth]{img/categ_corr_4w.png}
\small (b) Four wheeled vehicles
\end{minipage}
\caption{Cram\'er's $V$ association matrices for the two analytical populations.}
\label{fig:cramer-comparison}
\end{figure}
\FloatBarrier

\paragraph{Geography:} The \textbf{Bronx has the highest raw arrest rate in both populations}. For two wheeled stops, borough level rates range from approximately $3.2\%$ in Manhattan to $12.1\%$ in the Bronx. For four wheeled stops, we range from about $2.5\%$ in Queens to $4.8\%$ in the Bronx.

\paragraph{Hour:} It produces the most consistent raw pattern across the two populations. \textbf{Arrest rates are lowest during the morning and increase during late night and early morning hours}. Around 07:00 the raw rate is approximately $3.2\%$ for two wheeled vehicles and $0.6\%$ for four wheeled vehicles, while around 03:00 it rises to about $10.9\%$ and $5.3\%$. The shape is \textbf{clearly non linear} and motivates treating hour as more than a standard numerical predictor.

\paragraph{Sex:} \textbf{Male motorists have higher raw arrest rates in both datasets}. The difference is approximately $5.7\%$ versus $2.7\%$ for two wheeled stops and $3.70\%$ versus $1.51\%$ for four wheeled stops.

\paragraph{Ethnicity:} R\textbf{eported ethnicity shows the strongest demographic contrast}. In the two wheeled population, Hispanic and Black motorists have arrest rates of $7.24\%$ and $6.86\%$, compared with $3.23\%$ for White motorists. In the four wheeled population the corresponding rates are $4.66\%$, $5.02\%$ and $1.15\%$.

\paragraph{Age:} \textbf{Age provides the clearest example of disagreement between the populations}. Two wheeled arrest rates are highest among the youngest groups and do not follow the same ordering observed for cars. In the four wheeled population, the 25--34 and 18--24 groups have the highest rates, while arrest frequency declines strongly at older ages.

